\section*{4. 第 XVIII 讲的第二个补充}

本节在 $\mathbb C$ 上工作, 使用整系数经典上同调. 我们采用更接近 Lefschetz 的方法, 重新考察第 XVIII 讲的一些结果. 在那里, 这些结果只在与 $\mathbb Q$（更准确地说, 与 $\mathbb Q_\ell$, 但二者等价）张量后得到证明.

\subsubsection*{4.1.}
设 $X\subset\mathbb P$ 为纯维数 $n+1$ 的光滑复射影簇, $\ell$ 为 $X$ 的超平面截面的 Lefschetz 线束, 轴为 $A$. 也以 $\ell$ 表示线束的参数空间, 它同构于 $\mathbb P^1$. 设 $S\subset\ell$ 为例外值集. 对 $s\in\ell$, 以 $H_s$ 表示 $X$ 的相应超平面截面（维数为 $n$）.

设 $T(A)$ 为 $A\cap X$ 在 $X$ 中的一个闭管状邻域, 其边界与每个 $H_s$ 横截.

选择两点 $0,\infty\in\ell-S$, 并设 $\alpha:\ell\to\mathbb P^1(\mathbb C)$ 为保持定向的微分同胚, 将 $0$ 映到 $0$, $\infty$ 映到 $\infty$, 并将 $S$ 映到以 $0$ 为中心的单位圆周上.
\footnote{原印本将 $\alpha$ 的定义域写为 $\ell-S$, 但同一句中又求 $\alpha$ 在 $S$ 上的值; 所需定义域应为 $\ell$.}
$\alpha$ 与到 $\ell$ 的投影的复合 $\alpha f$, 将 $X-T(A)^\circ$ 变成 $\mathbb P^1(\mathbb C)$ 上以带边流形为纤维的奇异纤维化（$\alpha(S)$ 的每个点之上有一个奇点）. 以 $c_s$ 表示连接 $0$ 与 $\alpha(s)$ 的直线段.

在记号中省略 $\alpha$, 图景如下
\begin{center}
\begin{tikzpicture}[x=1.2cm,y=1.2cm]
  \coordinate (O) at (0,0);
  \coordinate (L) at (-1.8,1.8);
  \coordinate (S2) at (0,2.8);
  \coordinate (S1) at (2.2,2.0);
  \draw (O)--(L);
  \draw (O)--(S2);
  \draw (O)--(S1);
  \node[right] at (O) {$0$};
  \node[above right] at (S2) {$s_2$};
  \node[above right] at (S1) {$s_1$};
  \node[right] at (0,1.65) {$c_2$};
  \node[below right] at (1.25,1.14) {$c_1$};
  \fill (-0.73,0.76) circle (0.035);
  \fill (-0.55,0.91) circle (0.035);
  \fill (-0.37,1.03) circle (0.035);
\end{tikzpicture}
\end{center}

% Independent French verification: PDF346, printed336, XIX local9 says "z reel positif", not "z_i reel positif". The base coordinate z defines the path; z_i real defines the sphere fibers.
在 $s\in S$ 以及 $H_s\subset X$ 的奇点 $x_s$ 附近, 适当的局部坐标将我们置于 XIV 3.2 所研究的情形. 若选择 $\alpha$ 在 $s$ 附近全纯, 甚至可在 $X$ 和 $\ell$ 上选择局部坐标, 使 $f$ 写为 $\sum z_i^2$, 且 $\alpha^{-1}(c_s)$ 由“基底坐标 $z$ 为正实数”给出.\EditorialFootnote{法文原页（PDF 第 346 页，印页 336）此处是“$z$ 为正实数”，指基底坐标；英文工作版误作“各 $z_i$ 均为正实数”. 随后的球面丛条件仍是各 $z_i$ 为实数.} 在 $s$ 的一个邻域与 $\alpha^{-1}(c_s)-\{s\}$ 之交的上方, 考虑由“各 $z_i$ 均为实数”定义的实球面丛 $\Delta'_s\subset X$. 当 $t\to s$ 时, $\Delta'_s$ 在 $t$ 处的纤维 $\Delta'_{st}$ 趋于 $x_s$.

以 $\Delta_s$ 表示 $\{x_s\}$ 与 $\alpha^{-1}(c_s)-\{s\}$ 上的一个球面丛之并, 其中该球面丛包含于 $X-T(A)$ 并延拓 $\Delta'_s$. 则 $\Delta_s$ 是 $X$ 中的一个嵌入圆盘, 其边界 $H_0\cap\Delta_s$ 是一个球面, 代表沿路径 $\alpha^{-1}(c_s)$ 搬移到 $H_0$ 的 $s$ 处消灭圈 $\delta_s$.

对函数 $|\alpha f|^2$ 应用 Morse 理论, 研究
\[
 \{x\in X\mid x\in T(A)\ \text{或}\ |\alpha f(x)|^2\leq r\}
\]
的同伦型随 $r$ 增大如何变化. 置
\[
 X^- = \{x\in X\mid x\in T(A)\ \text{或}\ \alpha f(x)\ne\infty\}.
\]
可知包含映射
\[
 H_0\cup\bigcup_{s\in S}\Delta_s
 \hookrightarrow
 T(A)\cup H_0\cup\bigcup_{s\in S}\Delta_s
 \hookrightarrow X^-
\]
为同伦等价. Morse 理论还说明, $X-X^-$ 可形变收缩到一个实维数至多为 $n$ 的子复形上. 综合起来得到下述结果.

\subsubsection*{定理 4.2.}
$H_\infty^{\mathrm{aff}}=H_\infty-A$ 中存在实维数至多为 $n$ 的紧子复形 $\Sigma$, 使包含映射
\[
 H_0\cup\bigcup_{s\in S}\Delta_s\longrightarrow X-\Sigma
\]
为同伦等价.

\subsubsection*{4.3.}
去掉 $\Sigma$ 不改变 $X$ 在次数 $i\leq n$ 的上同调群、同调群或同伦群. 因为 $H_0\cup\bigcup_{s\in S}\Delta_s$ 是由 $H_0$ 沿代表消灭圈的球面黏附 $(n+1)$-维圆盘所得的, 同调给出\footnote{原印本在此列及下图中写作 $V_0$, 但第 4 节将纤维定义为 $H_0$.}
\[
 \mathbb Z^S\longrightarrow H_n(H_0)\longrightarrow H_n(X)
 \longrightarrow0.
\]
相应的上同调正合列, 连同由它经 Poincar\'e 对偶得到的列, 组成交换图
\begin{equation}\tag{4.3.1}
\begin{tikzcd}[row sep=large,column sep=large]
{} & \mathbb Z^S & {} & {} \\
\mathbb Z^S \arrow[r,"{\sum n_s\delta_s}"] &
H^n(H_0) \arrow[u,"{(\ ,\delta_s)}"'] \arrow[r] &
H^{n+2}(X) \arrow[r] & 0 \\
{} & H^n(X) \arrow[u] \arrow[ur,"{\eta\wedge}"'] & {} & {} \\
{} & 0 \arrow[u] & {} & {}
\end{tikzcd}
\end{equation}
